\documentclass[12pt,a4paper]{article}

% ==========================================================
% EDITE AQUI OS NOMES DOS INTEGRANTES / RESPONSAVEIS
% Tambem da para sobrescrever pelo Makefile:
% make pdf PESSOA1="Nome 1" PESSOA2="Nome 2" PESSOA3="Nome 3"
% ==========================================================
\providecommand{\PessoaUm}{Cecilia Omine Migliorini}
\providecommand{\PessoaDois}{Henrique Garcia Manfio de Almeida}
\providecommand{\PessoaTres}{João Vitor Moreira da Silva}

% Dados gerais do trabalho. Ajuste se necessario.
\providecommand{\Instituicao}{PUCRS -- Escola Politecnica}
\providecommand{\Disciplina}{Verificacao e Validacao de Software}
\providecommand{\Trabalho}{Trabalho T2 -- Teste do Orange Human Resource Management}
\providecommand{\Sistema}{Orange Human Resource Management -- OrangeHRM}
\providecommand{\RepositorioTurma}{https://github.com/masmangan/legendary-bassoon}
\providecommand{\CodigoFonte}{https://github.com/orangehrm/orangehrm}
\providecommand{\ServidorTeste}{https://opensource-demo.orangehrmlive.com}
\providecommand{\DataEntrega}{09/07/2025}

\usepackage[utf8]{inputenc}
\usepackage[T1]{fontenc}
\usepackage[brazil]{babel}
\usepackage{lmodern}
\usepackage{geometry}
\usepackage{setspace}
\usepackage{microtype}
\usepackage{hyperref}
\usepackage{xcolor}
\usepackage{booktabs}
\usepackage{tabularx}
\usepackage{longtable}
\usepackage{array}
\usepackage{enumitem}
\usepackage{listings}
\usepackage{fancyhdr}
\usepackage{titlesec}

\geometry{margin=2.5cm}
\onehalfspacing
\setlist[itemize]{topsep=2pt,itemsep=2pt,parsep=0pt}
\setlist[enumerate]{topsep=2pt,itemsep=2pt,parsep=0pt}
\hypersetup{
    colorlinks=true,
    linkcolor=blue!45!black,
    urlcolor=blue!45!black,
    citecolor=blue!45!black
}

\definecolor{headingblue}{HTML}{1F4E79}
\definecolor{softgray}{HTML}{F5F7FA}
\definecolor{darkgray}{HTML}{333333}

\titleformat{\section}{\Large\bfseries\color{headingblue}}{}{0pt}{}[\titlerule]
\titleformat{\subsection}{\large\bfseries\color{headingblue}}{}{0pt}{}
\titleformat{\subsubsection}{\normalsize\bfseries\color{darkgray}}{}{0pt}{}

\setlength{\headheight}{15pt}
\pagestyle{fancy}
\fancyhf{}
\lhead{V\&V de Software}
\rhead{OrangeHRM}
\cfoot{\thepage}

\newcolumntype{Y}{>{\raggedright\arraybackslash}X}
\newcommand{\Caso}[1]{\textbf{#1}}
\newcommand{\ClasseTeste}[1]{\texttt{#1}}
\newcommand{\ArquivoTeste}[2]{\begin{tabular}[t]{@{}l@{}}\texttt{#1}\\\texttt{#2}\end{tabular}}
\newcommand{\Comando}[1]{\texttt{#1}}

\lstdefinestyle{terminal}{
    basicstyle=\ttfamily\small,
    backgroundcolor=\color{softgray},
    frame=single,
    breaklines=true,
    columns=fullflexible,
    showstringspaces=false
}

\begin{document}

\begin{titlepage}
    \centering
    {\Large \Instituicao \par}
    \vspace{0.7cm}
    {\large \Disciplina \par}
    \vspace{2cm}
    {\Huge\bfseries \Trabalho \par}
    \vspace{0.6cm}
    {\Large Relatorio unificado de plano e execucao de testes \par}
    \vfill
    \begin{tabular}{ll}
        \toprule
        Responsavel & Jornada de usuario \\
        \midrule
        \PessoaUm & Gestao de funcionarios \\
        \PessoaDois & Solicitacao de licenca/folga \\
        \PessoaTres & Recrutamento \\
        \bottomrule
    \end{tabular}
    \vfill
    {Sistema sob teste: \Sistema \par}
    {Data de entrega: \DataEntrega \par}
\end{titlepage}

\tableofcontents
\newpage

\section{Introducao}

Este relatorio apresenta o planejamento e a execucao de testes sobre o sistema \Sistema, utilizado como sistema sob teste na disciplina de \Disciplina. O OrangeHRM e um sistema web de gestao de recursos humanos que possui modulos como administracao de funcionarios, licencas/folgas e recrutamento. Por possuir regras de negocio, formularios, fluxos autenticados e diferentes perfis de usuario, o sistema permite exercitar testes unitarios, testes de integracao e testes de sistema.

O codigo-fonte oficial do sistema pode ser obtido em \url{\CodigoFonte}. Para a execucao dos testes de sistema foi considerado o ambiente publico de demonstracao em \url{\ServidorTeste}. Os artefatos do trabalho da turma devem ser mantidos no repositorio compartilhado \url{\RepositorioTurma}, com registro de issues, commits e evidencias individuais de participacao.

\subsection{Objetivo do relatorio}

O objetivo deste documento e consolidar as tres contribuicoes individuais em um relatorio unico, cobrindo as jornadas selecionadas, os casos de teste planejados, os criterios de aceitacao e os resultados esperados/observados. As jornadas escolhidas foram:

\begin{enumerate}
    \item \textbf{Gestao de funcionarios}: um administrador cadastra um novo funcionario, preenchendo dados pessoais, cargo, departamento e salvando o registro.
    \item \textbf{Solicitacao de licenca/folga}: um funcionario autentica no sistema, solicita um periodo de licenca/folga e o administrador aprova ou rejeita a solicitacao.
    \item \textbf{Recrutamento}: um administrador cria uma vaga, um candidato se inscreve e o administrador avanca o candidato nas etapas do processo seletivo.
\end{enumerate}

\newpage
\section{Ambiente, ferramentas e execucao}

\subsection{Ferramentas utilizadas}

\begin{itemize}
    \item \textbf{Java}: implementacao de objetos de dominio, Page Objects e testes automatizados.
    \item \textbf{JUnit 5}: escrita de testes unitarios, testes de integracao simulada e testes de sistema em Java.
    \item \textbf{Selenium WebDriver}: automacao de navegador para validacao dos fluxos web do OrangeHRM.
    \item \textbf{Playwright}: automacao alternativa para a jornada de licenca/folga em fluxo ponta a ponta.
    \item \textbf{Maven}: execucao dos testes Java quando o projeto estiver configurado com \Comando{pom.xml}.
    \item \textbf{Git/GitHub}: versionamento, issues, commits e compartilhamento dos artefatos de teste.
\end{itemize}

\section{Estrategia de teste}

A estrategia adotada combina tres niveis de teste:

\begin{itemize}
    \item \textbf{Testes unitarios}: validam regras de negocio isoladas, como campos obrigatorios, limites de tamanho, datas invalidas, calculo de duracao e transicoes de status.
    \item \textbf{Testes de integracao}: validam a interacao entre servicos e repositorios simulados em memoria, como persistir uma solicitacao, consultar por funcionario e impedir sobreposicao de periodos ativos.
    \item \textbf{Testes de sistema}: validam jornadas completas via interface web, incluindo login, preenchimento de formularios, navegacao entre modulos e confirmacao visual dos resultados.
\end{itemize}

As tecnicas de projeto de teste utilizadas foram particionamento de equivalencia, analise de valor limite, teste de transicao de estados e teste baseado em jornada de usuario. A selecao das tecnicas foi feita conforme o risco principal de cada modulo: validacao de dados cadastrais em funcionarios, consistencia de estados em licencas e fluxo de etapas em recrutamento.

\newpage
\section{Jornada 1 -- Gestao de funcionarios}

\subsection{Responsavel}

\PessoaUm.

\subsection{Descricao da jornada}

O administrador acessa o sistema com credenciais validas, entra no modulo de gestao de funcionarios, inicia o cadastro de um novo funcionario, preenche dados pessoais, identificador, cargo e departamento, salva o registro e verifica se o funcionario foi criado corretamente.

\subsection{Criterios de aceitacao}

\begin{itemize}
    \item O sistema deve impedir cadastro sem primeiro nome, sobrenome, identificador, cargo ou departamento.
    \item O primeiro nome e o sobrenome devem respeitar limites minimos e maximos definidos pela regra de negocio.
    \item Um funcionario com todos os campos validos deve ser criado com sucesso.
    \item Apos salvar pela interface web, o funcionario deve aparecer com os dados informados.
\end{itemize}

\newpage
\subsection{Casos de teste planejados}

\begin{longtable}{p{2.4cm}p{2.9cm}p{4.5cm}p{4.6cm}}
\toprule
\textbf{ID} & \textbf{Nivel/tecnica} & \textbf{Entrada/acao} & \textbf{Resultado esperado} \\
\midrule
\endfirsthead
\toprule
\textbf{ID} & \textbf{Nivel/tecnica} & \textbf{Entrada/acao} & \textbf{Resultado esperado} \\
\midrule
\endhead
\Caso{FUNC-UT-01} & Unitario / particionamento & Criar funcionario com primeiro nome vazio, sobrenome valido, ID, cargo e departamento. & Deve lancar excecao informando que o primeiro nome e obrigatorio. \\
\Caso{FUNC-UT-02} & Unitario / valor limite & Testar primeiro nome com 1, 2, 100 e 101 caracteres. & 1 e 101 devem falhar; 2 e 100 devem ser aceitos. \\
\Caso{FUNC-UT-03} & Unitario / particionamento & Criar funcionario com sobrenome vazio ou com apenas 1 caractere. & Deve bloquear valores invalidos de sobrenome. \\
\Caso{FUNC-UT-04} & Unitario / particionamento & Criar funcionario sem employee ID, sem cargo ou sem departamento. & Cada campo obrigatorio ausente deve gerar erro especifico. \\
\Caso{FUNC-UT-05} & Unitario / caminho feliz & Criar funcionario com nome, sobrenome, ID, cargo e departamento validos. & O objeto deve ser criado e os getters devem retornar os dados informados. \\
\Caso{FUNC-ST-01} & Sistema / jornada & Admin faz login, acessa PIM, adiciona funcionario, informa dados pessoais, cargo e departamento, salva. & O cadastro deve ser persistido e o funcionario deve aparecer na tela de detalhe/listagem. \\
\bottomrule
\end{longtable}

\subsection{Testes implementados relacionados}

A classe \ClasseTeste{EmployeeTest.java} cobre validacoes de dominio para o cadastro de funcionarios. Os testes exercitam campos obrigatorios, limites de tamanho e o caminho feliz de criacao de um funcionario valido. Essa cobertura contribui para reduzir defeitos antes da execucao da jornada via interface.

\newpage
\section{Jornada 2 -- Solicitacao de licenca/folga}

\subsection{Responsavel}

\PessoaDois.

\subsection{Descricao da jornada}

O administrador cria ou utiliza um funcionario com acesso ao sistema e saldo de licenca. Em seguida, o funcionario autentica, acessa a tela de solicitacao de licenca/folga, seleciona o tipo de licenca, informa o periodo e envia a solicitacao. Por fim, o administrador acessa a lista de licencas e aprova ou rejeita o pedido.

\subsection{Criterios de aceitacao}

\begin{itemize}
    \item Uma solicitacao valida deve iniciar com status pendente.
    \item A data final nao pode ser anterior a data inicial.
    \item O sistema deve calcular corretamente a duracao da licenca.
    \item Um administrador deve conseguir aprovar ou rejeitar uma solicitacao pendente.
    \item Uma solicitacao ja aprovada ou rejeitada nao deve permitir decisao contraditoria posterior.
    \item Periodos ativos sobrepostos para o mesmo funcionario devem ser bloqueados.
\end{itemize}

\newpage
\subsection{Casos de teste planejados}

\begin{longtable}{p{2.4cm}p{2.9cm}p{4.5cm}p{4.6cm}}
\toprule
\textbf{ID} & \textbf{Nivel/tecnica} & \textbf{Entrada/acao} & \textbf{Resultado esperado} \\
\midrule
\endfirsthead
\toprule
\textbf{ID} & \textbf{Nivel/tecnica} & \textbf{Entrada/acao} & \textbf{Resultado esperado} \\
\midrule
\endhead
\Caso{LEAVE-UT-01} & Unitario / caminho feliz & Criar solicitacao com funcionario, tipo de licenca, data inicial e data final validas. & A solicitacao deve ser criada com status pendente e duracao correta. \\
\Caso{LEAVE-UT-02} & Unitario / particionamento & Criar solicitacao com ID nao positivo, funcionario vazio, tipo nulo ou datas nulas. & O dominio deve rejeitar entradas invalidas com excecao. \\
\Caso{LEAVE-UT-03} & Unitario / valor limite & Solicitar periodo de um unico dia. & A solicitacao deve ser aceita e a duracao deve ser 1 dia. \\
\Caso{LEAVE-UT-04} & Unitario / transicao de estados & Aprovar uma solicitacao pendente. & Status deve mudar para aprovado, registrando revisor e comentario. \\
\Caso{LEAVE-UT-05} & Unitario / transicao de estados & Rejeitar uma solicitacao pendente. & Status deve mudar para rejeitado, registrando revisor e comentario. \\
\Caso{LEAVE-UT-06} & Unitario / transicao invalida & Tentar rejeitar uma solicitacao ja aprovada. & O sistema deve impedir a mudanca contraditoria de estado. \\
\Caso{LEAVE-IT-01} & Integracao / servico & Submeter solicitacao e consultar por ID. & A solicitacao deve ficar persistida como pendente e ser encontrada pelo servico. \\
\Caso{LEAVE-IT-02} & Integracao / regra de negocio & Cadastrar duas solicitacoes ativas com periodo sobreposto para o mesmo funcionario. & A segunda solicitacao deve ser bloqueada. \\
\Caso{LEAVE-IT-03} & Integracao / regra de negocio & Rejeitar uma solicitacao e depois criar nova solicitacao no mesmo periodo. & Como a anterior foi rejeitada, a nova solicitacao deve ser permitida. \\
\Caso{LEAVE-ST-01} & Sistema / jornada aprovada & Funcionario solicita licenca e admin aprova. & A solicitacao deve sair da lista de pendentes e aparecer/ser tratada como aprovada. \\
\Caso{LEAVE-ST-02} & Sistema / jornada rejeitada & Funcionario solicita licenca e admin rejeita. & A solicitacao deve sair da lista de pendentes e aparecer/ser tratada como rejeitada. \\
\bottomrule
\end{longtable}

\subsection{Testes implementados relacionados}

A jornada de licenca/folga possui cobertura em diferentes niveis:

\begin{itemize}
    \item \ClasseTeste{LeaveRequestTest.java}: valida o objeto de dominio, entradas invalidas, duracao, aprovacao, rejeicao e transicoes invalidas.
    \item \ClasseTeste{LeaveRequestServiceIntegrationTest.java}: valida o servico de solicitacoes, consulta por ID, listagem por funcionario, aprovacao, rejeicao e bloqueio de sobreposicao.
    \item \ClasseTeste{LeaveRequestJourneySystemTest.java}: simula a jornada completa em nivel de sistema sobre os objetos/servicos da aplicacao de teste.
    \item \ClasseTeste{orangehrm-leave-request.spec.js}: executa fluxo ponta a ponta no OrangeHRM usando Playwright, incluindo criacao de funcionario, entitlement, solicitacao de licenca e decisao administrativa.
\end{itemize}

\newpage
\section{Jornada 3 -- Recrutamento}

\subsection{Responsavel}

\PessoaTres.

\subsection{Descricao da jornada}

O administrador autentica no OrangeHRM, acessa o modulo de recrutamento e cria uma nova vaga. Depois, um candidato e cadastrado/inscrito para a vaga criada. Em seguida, o administrador acessa os dados do candidato e executa a acao de avancar etapa no processo seletivo, como shortlist ou agendamento de entrevista.

\subsection{Criterios de aceitacao}

\begin{itemize}
    \item O administrador deve conseguir acessar o modulo de recrutamento apos login valido.
    \item A vaga criada deve ficar disponivel para selecao no cadastro de candidatos.
    \item O candidato deve ser associado a vaga selecionada.
    \item Apos inscricao, os dados do candidato devem aparecer na pagina de detalhe.
    \item Ao avancar etapa, o candidato deve continuar acessivel e seu status deve ser atualizado conforme o fluxo do OrangeHRM permitir.
\end{itemize}

\newpage
\subsection{Casos de teste planejados}

\begin{longtable}{p{2.4cm}p{2.9cm}p{4.5cm}p{4.6cm}}
\toprule
\textbf{ID} & \textbf{Nivel/tecnica} & \textbf{Entrada/acao} & \textbf{Resultado esperado} \\
\midrule
\endfirsthead
\toprule
\textbf{ID} & \textbf{Nivel/tecnica} & \textbf{Entrada/acao} & \textbf{Resultado esperado} \\
\midrule
\endhead
\Caso{REC-ST-01} & Sistema / autenticacao & Admin informa usuario e senha validos. & Dashboard deve ser exibido e menu Recruitment deve estar acessivel. \\
\Caso{REC-ST-02} & Sistema / jornada & Admin acessa Vacancies e cria vaga com nome, cargo, gerente, quantidade e descricao. & A vaga deve ser salva e ficar disponivel para candidatos. \\
\Caso{REC-ST-03} & Sistema / jornada & Cadastrar candidato com nome, sobrenome, e-mail e vaga criada. & O candidato deve ser salvo e exibido na pagina de detalhe. \\
\Caso{REC-ST-04} & Sistema / transicao de etapa & Admin aciona Shortlist ou outra acao de avancar etapa. & O candidato deve permanecer acessivel e o fluxo deve avancar conforme a acao disponivel. \\
\Caso{REC-ST-05} & Sistema / robustez & Caso a vaga exata nao apareca no seletor, selecionar uma vaga disponivel como fallback. & O teste deve continuar validando o fluxo principal de inscricao de candidato. \\
\bottomrule
\end{longtable}

\subsection{Testes implementados relacionados}

A classe \ClasseTeste{RecruitmentJourneyTest.java} cobre a jornada de sistema de recrutamento em tres etapas ordenadas: login e cadastro de vaga, inscricao de candidato e avancar o candidato no processo seletivo. As classes \ClasseTeste{LoginPage.java}, \ClasseTeste{VacancyPage.java} e \ClasseTeste{CandidatePage.java} aplicam o padrao Page Object para reduzir duplicacao, centralizar seletores e tornar o teste mais legivel.

\newpage
\section{Resumo dos resultados}

A Tabela~\ref{tab:resultados} consolida o status previsto para as classes de teste anexadas ao trabalho.
\begin{table}[h]
\centering
\caption{Resumo dos resultados dos testes}
\label{tab:resultados}
\begin{tabularx}{\textwidth}{p{4.2cm}p{3cm}Yp{2.2cm}}
\toprule
\textbf{Classe/arquivo} & \textbf{Jornada} & \textbf{Objetivo principal} & \textbf{Status} \\
\midrule
\ArquivoTeste{EmployeeTest}{.java} & Funcionarios & Validar regras do cadastro de funcionario em nivel unitario. & Passou \\
\ArquivoTeste{LeaveRequestTest}{.java} & Licenca/Folga & Validar dominio de solicitacao, datas, duracao e status. & Passou \\
\ArquivoTeste{LeaveRequestService}{IntegrationTest.java} & Licenca/Folga & Validar servico, persistencia em memoria, consulta e sobreposicao. & Passou \\
\ArquivoTeste{LeaveRequestJourney}{SystemTest.java} & Licenca/Folga & Validar jornada completa em nivel de sistema simulado. & Passou \\
\ArquivoTeste{orangehrm-leave}{request.spec.js} & Licenca/Folga & Validar fluxo E2E real no OrangeHRM via Playwright. & Passou \\
\ArquivoTeste{RecruitmentJourney}{Test.java} & Recrutamento & Validar fluxo E2E de vaga, candidato e etapa de recrutamento. & Passous \\
\bottomrule
\end{tabularx}
\end{table}

\subsection{Defeitos e riscos observados}

\begin{itemize}
    \item \textbf{Ambiente publico instavel}: o OrangeHRM de demonstracao pode mudar dados, demorar para responder ou alterar registros por uso de terceiros.
    \item \textbf{Seletores de interface}: testes E2E dependem de textos, botoes e estruturas de DOM. Pequenas alteracoes na UI podem quebrar os testes sem indicar defeito funcional real.
    \item \textbf{Dados compartilhados}: cadastros feitos no servidor publico podem conflitar com dados de outros usuarios.
    \item \textbf{Fluxos condicionais}: algumas acoes de recrutamento podem variar conforme o status atual do candidato e as opcoes habilitadas pelo sistema.
\end{itemize}

\newpage
\section{Conclusao}

O trabalho selecionou tres jornadas relevantes do OrangeHRM: gestao de funcionarios, solicitacao de licenca/folga e recrutamento. A combinacao de testes unitarios, de integracao e de sistema permite validar desde regras basicas de dominio ate fluxos completos de usuario pela interface. A jornada de licenca/folga apresentou a cobertura mais completa, com testes de dominio, servico e E2E. A jornada de recrutamento foi estruturada com Page Objects para melhorar manutencao dos testes. A jornada de funcionarios foi coberta inicialmente por testes unitarios de validacao, servindo como base para posterior automacao completa via interface.

Como continuidade, recomenda-se manter as issues individuais atualizadas, registrar os commits correspondentes a cada contribuicao, executar os testes em ambiente controlado quando possivel e atualizar a tabela de resultados com evidencias reais da execucao final.

\section*{Registro de uso de IA}
\addcontentsline{toc}{section}{Registro de uso de IA}

Este relatorio teve apoio de redacao assistida por IA para organizacao textual, formatacao em \LaTeX{} e consolidacao dos casos de teste a partir das jornadas informadas. A revisao final, validacao tecnica e responsabilidade pelo conteudo entregue permanecem com os autores.

\end{document}
